%! Polynomial Function Behavior — Unit 3, Lesson 3.2 — Slides (Beamer)
% Phase order: hook -> I Do -> We Do -> You Do -> debrief, closing on the DeltaMath assignment.
% Frame labels name the notes rows and problem numbers (1–18) so the deck and the packet agree.
\documentclass[aspectratio=169, 11pt]{beamer}
\usepackage{precalculus-beamer}   % provides \plumheader{} and \sectionlabel[color]{}

\begin{document}

% TITLE SLIDE
{
\setbeamercolor{background canvas}{bg=plum}
\begin{frame}[plain]
  \vspace{2.8cm}\hspace{0.55cm}%
  \begin{minipage}{0.9\textwidth}
    {\color{goldacc}\bfseries\small LESSON 3.2}\\[0.5cm]
    {\color{white}\bfseries\LARGE Polynomial Function Behavior}\\[1.2cm]
    {\color{lilacmid}\small Precalculus~$\cdot$~Unit 3: Polynomial Functions}
  \end{minipage}
\end{frame}
}

% HOOK SLIDE
\begin{frame}[t, plain]
  \plumheader{Hook: A Race Between Two Functions}
  \sectionlabel[goldacc]{Think About It}

  \begin{center}
    {\Large\bfseries\color{plum} \(f(x) = x^4 \qquad\text{vs.}\qquad g(x) = 500x^3\)}
  \end{center}

  \medskip
  \begin{columns}[T]
    \column{0.48\textwidth}
      \begin{block}{Who is ahead?}
        At \(x = 10\)? \\[3pt]
        At \(x = 1000\)?
      \end{block}

    \column{0.48\textwidth}
      \begin{block}{Then ask}
        Who wins in the \textbf{long run}? Does the \textbf{500} decide it --- or the
        \textbf{exponent}?
      \end{block}
  \end{columns}

  \bigskip
  \textit{Far from zero, the highest power always takes over. Today that term gets a name.}
\end{frame}

% ROW 1 — ANATOMY
\begin{frame}[t, plain]
  \plumheader{Anatomy of a Polynomial}
  \sectionlabel[goldacc]{Notes row ANATOMY $\cdot$ I Do: problem 1 $\cdot$ We Do: problem 2}

  \begin{block}{Standard form, degree \(n\)}
    \(p(x) = a_n x^n + a_{n-1}x^{n-1} + \dots + a_1 x + a_0\), \quad \(n\) a whole number,
    \(a_n \ne 0\)
  \end{block}

  \medskip
  \begin{columns}[T]
    \column{0.5\textwidth}
      \begin{center}
        {\Large \(p(x) = -3x^4 + 2x^3 - x + 7\)}
      \end{center}
      \begin{itemize}\setlength{\itemsep}{4pt}
        \item \textbf{Degree:} 4 \quad (the highest power)
        \item \textbf{Leading term:} \(-3x^4\)
        \item \textbf{Leading coefficient:} \(-3\) \quad (sign included)
        \item \textbf{Constant term:} 7 \quad (\(= p(0)\))
      \end{itemize}

    \column{0.46\textwidth}
      \textbf{Try it.}
      \begin{enumerate}\setlength{\itemsep}{6pt}
        \item \(f(x) = 4x^3 - x^2 + 6\)
        \item \(g(x) = \tfrac12 x^6 - x\) \quad and \quad \(h(x) = 9\)
      \end{enumerate}
      \smallskip
      A nonzero constant has degree \textbf{0}.
  \end{columns}
\end{frame}

% ROW 2 — END BEHAVIOR
\begin{frame}[t, plain]
  \plumheader{End Behavior: Degree and Sign}
  \sectionlabel[goldacc]{Notes row END BEHAVIOR $\cdot$ I Do: problem 3 $\cdot$ We Do: problems 4--6}

  \begin{columns}[T]
    \column{0.52\textwidth}
      \begin{center}
      \renewcommand{\arraystretch}{1.2}
      \begin{tabular}{c|c|c}
       & \textbf{\(n\) even} & \textbf{\(n\) odd} \\ \hline
      \(a_n > 0\) &
      \begin{tikzpicture}[x=0.32cm, y=0.12cm]
        \draw[linegray] (-2.4,0) -- (2.4,0);
        \draw[very thick, plum, <->, samples=40, domain=-2.04:2.04] plot (\x, {\x*\x*\x*\x-3*\x*\x});
      \end{tikzpicture} &
      \begin{tikzpicture}[x=0.32cm, y=0.12cm]
        \draw[linegray] (-2.4,0) -- (2.4,0);
        \draw[very thick, plum, <->, samples=40, domain=-2.2:2.2] plot (\x, {\x*\x*\x-3*\x});
      \end{tikzpicture} \\
       & \scriptsize up, up & \scriptsize down, up \\ \hline
      \(a_n < 0\) &
      \begin{tikzpicture}[x=0.32cm, y=0.12cm]
        \draw[linegray] (-2.4,0) -- (2.4,0);
        \draw[very thick, plum, <->, samples=40, domain=-2.04:2.04] plot (\x, {-\x*\x*\x*\x+3*\x*\x});
      \end{tikzpicture} &
      \begin{tikzpicture}[x=0.32cm, y=0.12cm]
        \draw[linegray] (-2.4,0) -- (2.4,0);
        \draw[very thick, plum, <->, samples=40, domain=-2.2:2.2] plot (\x, {-\x*\x*\x+3*\x});
      \end{tikzpicture} \\
       & \scriptsize down, down & \scriptsize up, down \\
      \end{tabular}
      \end{center}

    \column{0.44\textwidth}
      \begin{block}{The leading term decides}
        Even degree: the arms \textbf{match}. Odd: they \textbf{oppose}. A positive
        leading coefficient sends the \textbf{right} arm up.
      \end{block}
      \smallskip
      \(p(x) = 2x^4 - 5x^3 + x\):
      \[ \lim_{x \to -\infty} p(x) = \infty, \qquad \lim_{x \to \infty} p(x) = \infty \]
      \textbf{Even degree} \(\Rightarrow\) an absolute max (arms down) or min (arms up).
  \end{columns}
\end{frame}

% ROW 3 — COUNTING
\begin{frame}[t, plain]
  \plumheader{At Most, Not Exactly}
  \sectionlabel[goldacc]{Notes row COUNTING $\cdot$ I Do: problem 7 $\cdot$ We Do: problems 8--10}

  \begin{block}{Two ceilings for degree \(n\)}
    At most \(\mathbf{n - 1}\) turning points. \qquad At most \(\mathbf{n}\) real zeros
    (\(x\)-intercepts).
  \end{block}

  \medskip
  \begin{columns}[T]
    \column{0.3\textwidth}
      \begin{center}
      \begin{tikzpicture}[x=0.36cm, y=0.18cm]
        \draw[linegray] (-2.8,0) -- (2.8,0); \draw[linegray] (0,-5) -- (0,6);
        \draw[very thick, plum, <->, samples=50, domain=-2.5:2.5] plot (\x, {0.75*(\x*\x*\x-4*\x)});
      \end{tikzpicture}\\
      \small A: degree 3 --- 2 turns, 3 zeros
      \end{center}
    \column{0.3\textwidth}
      \begin{center}
      \begin{tikzpicture}[x=0.36cm, y=0.18cm]
        \draw[linegray] (-2.8,0) -- (2.8,0); \draw[linegray] (0,-5) -- (0,6);
        \draw[very thick, plum, <->, samples=50, domain=-2.25:2.25] plot (\x, {0.5*\x*\x*\x*\x-2*\x*\x+3});
      \end{tikzpicture}\\
      \small B: degree 4 --- 3 turns, \textbf{0 zeros}
      \end{center}
    \column{0.3\textwidth}
      \begin{center}
      \begin{tikzpicture}[x=0.36cm, y=0.18cm]
        \draw[linegray] (-2.8,0) -- (2.8,0); \draw[linegray] (0,-5) -- (0,6);
        \draw[very thick, plum, <->, samples=50, domain=-2.05:2.3] plot (\x, {1-0.5*\x*\x*\x});
      \end{tikzpicture}\\
      \small C: degree 3 --- \textbf{0 turns}, 1 zero
      \end{center}
  \end{columns}

  \medskip
  \textbf{Read backwards:} \(k\) turning points \(\Rightarrow\) degree \(\ge k + 1\); \quad
  \(m\) \(x\)-intercepts \(\Rightarrow\) degree \(\ge m\).
\end{frame}

% ROW 4 — THE CRUX
\begin{frame}[t, plain]
  \plumheader{Which Term Leads?}
  \sectionlabel[redacc]{Notes row WHICH TERM LEADS? $\cdot$ I Do: problem 11 $\cdot$ We Do: 12--14 (13 is the crux)}

  \begin{columns}[T]
    \column{0.46\textwidth}
      \textbf{Out of order}
      \[ p(x) = 5 + 3x - 2x^4 \]
      Written first: \(5\). \quad Leading term: \(\mathbf{-2x^4}\).

      \smallskip
      Even, negative: \(\displaystyle\lim_{x \to \pm\infty} p(x) = -\infty\).

    \column{0.50\textwidth}
      \textbf{Factored form --- never expand}
      \[ q(x) = -(x-1)^2(x+3)(2x+1) \]
      Leading pieces: \((-1)(x^2)(x)(2x) = \mathbf{-2x^4}\)

      \smallskip
      Degree \(2 + 1 + 1 = 4\); leading coefficient \(-2\). Arms down: an absolute max.
  \end{columns}

  \bigskip
  \begin{block}{The rule}
    The leading term is the term with the \textbf{highest power}, wherever it is written.
    Find the highest power first; then read its sign.
  \end{block}

  \smallskip
  \textbf{Now:} \(f(x) = 3x^2 - x^5 + 1\) \quad and \quad \(s(x) = (4 - x)(x + 2)^2\).
\end{frame}

% YOU DO
\begin{frame}[t, plain]
  \plumheader{You Do --- Read It All}
  \sectionlabel[redacc]{Notes row READ IT ALL $\cdot$ problems 15--18, alone, 16 minutes}

  \begin{columns}[T]
    \column{0.62\textwidth}
      \begin{enumerate}\setlength{\itemsep}{7pt}
        \setcounter{enumi}{14}
        \item \(f(x) = 6x - x^3 + 2x^2 - 8\). Degree, leading coefficient, end behavior,
              the most turning points, the most real zeros.
        \item \(g(x) = (2x - 1)^3(5 - x)\). Degree, leading coefficient, end behavior.
              Absolute max or min?
        \item Graph D: end behavior, sign of the leading coefficient, least degree. Could
              it be degree 5? Degree 6?
        \item Can a polynomial of degree 3 have \textbf{no} real zeros? Use end behavior.
      \end{enumerate}
    \column{0.34\textwidth}
      \begin{center}
      \begin{tikzpicture}[x=0.5cm, y=0.28cm]
        \draw[linegray, very thin, step=1] (-3,-5) grid (3,3);
        \draw[thick] (-3,0) -- (3,0); \draw[thick] (0,-5) -- (0,3);
        \draw[very thick, plum, <->, samples=60, domain=-2.3:2.3] plot (\x, {-0.8*(\x*\x-1)*(\x*\x-4)});
      \end{tikzpicture}\\[2pt]
      \small Graph D
      \end{center}
  \end{columns}
\end{frame}

% DEBRIEF
\begin{frame}[t, plain]
  \plumheader{Debrief --- The Headline}
  \sectionlabel[goldacc]{Share out 15--18, one answer each $\cdot$ then say it back}

  \begin{columns}[T]
    \column{0.48\textwidth}
      \begin{block}{The leading term}
        The \textbf{highest power}, wherever it is written.\\[3pt]
        Its \textbf{degree} and its \textbf{sign} decide both ends.
      \end{block}
    \column{0.48\textwidth}
      \begin{block}{The ceilings}
        At most \(n - 1\) turning points.\\[3pt]
        At most \(n\) real zeros --- \textbf{at most}, not exactly.
      \end{block}
  \end{columns}

  \bigskip
  \textbf{After the break (3.3):} the zeros themselves --- where the graph crosses, where it
  only touches, and why the ceiling of \(n\) becomes an exact count once complex zeros are
  counted.
\end{frame}

% HOMEWORK — DELTAMATH
\begin{frame}[t, plain]
  \plumheader{Homework: DeltaMath}
  \sectionlabel[redacc]{Graded $\cdot$ online $\cdot$ assignment \textit{Lesson 3.2}}

  \begin{itemize}\setlength{\itemsep}{8pt}
    \item Your graded homework is the DeltaMath assignment \textbf{Lesson 3.2}.
          Copy the due date into the \textbf{Homework} row of your cover sheet:
          \quad Due: \rule{3cm}{0.4pt}
    \item The topics, in order:
          \begin{itemize}\setlength{\itemsep}{2pt}
            \item Polynomial terminology
            \item End behavior --- from the formula (find the leading term \textit{first}),
                  in limit notation, and from a graph
            \item Even and odd degree from a graph
            \item Features of a polynomial graph --- turning points and intercepts
          \end{itemize}
    \item \textbf{AP Practice} (the last two pages of the packet) is
          \textbf{extra credit, scored 1 to 5}. Optional, AP-exam style. Do DeltaMath first.
  \end{itemize}
\end{frame}

\end{document}
